\documentclass[11pt,a4paper]{article}

\usepackage[utf8]{inputenc}
\usepackage[T1]{fontenc}
\usepackage{amsmath,amssymb}
\usepackage{booktabs}
\usepackage{graphicx}
\usepackage{hyperref}
\usepackage[margin=2.5cm]{geometry}
\usepackage{algorithm}
\usepackage{algorithmic}
\usepackage{multirow}

\title{Split Validation: Structural Regularization for\\Joint Probability Trees via Asymmetric Data Partitioning}

\author{Technical Report}
\date{March 2026}

\begin{document}
\maketitle

\begin{abstract}
We introduce \emph{split validation}, a structural regularization technique for
Joint Probability Tree (JPT) learning that partitions the training data into
disjoint training and evaluation subsets at the split-selection level. During
tree induction, only training samples provide candidate split points, while
target values from all (or a configurable subset of) samples contribute to the
impurity score. This asymmetry forces each split to generalize beyond training
feature noise. We evaluate the method on synthetic classification and regression
tasks as well as three real-world datasets (Iris, Wine, Breast Cancer) using
accuracy, mean squared error, and log-likelihood metrics. Results show that
split validation consistently reduces tree complexity (6\% fewer leaves,
$p < 0.001$), significantly improves test log-likelihood ($p = 0.023$), and
yields notable accuracy gains on real-world data---particularly on the Wine
dataset, where a 0.3 training ratio improves accuracy from 0.877 to 0.939.
The method introduces zero computational overhead when disabled and minimal
overhead when active.
\end{abstract}

% =========================================================================
\section{Introduction}
\label{sec:intro}

Joint Probability Trees~(JPTs) are density estimation models that combine the
interpretability of decision trees with the expressiveness of local
distribution estimation. Like all tree-based learners, JPTs are susceptible to
overfitting when the tree grows too deep, splitting on noise rather than signal
in the feature values.

Standard regularization techniques for decision trees include pre-pruning
(minimum samples per leaf, maximum depth, minimum impurity improvement) and
post-pruning (reduced-error pruning, cost-complexity pruning). These approaches
either restrict tree growth globally or evaluate subtrees after construction.

We propose \emph{split validation}, which operates at a different granularity:
within each individual split decision. The key idea is simple:

\begin{enumerate}
    \item Partition the training data into \emph{training} and \emph{evaluation}
          samples via a boolean mask.
    \item When searching for the best split on a feature, only \emph{training}
          sample values serve as candidate split points.
    \item The impurity score for each candidate split is computed using target
          values from \emph{all} samples (or a configurable subset).
\end{enumerate}

This creates an asymmetry: a split must explain both training and evaluation
targets well, but can only place the split boundary at positions defined by
training feature values. If a training feature value is an outlier that happens
to separate training targets well but not evaluation targets, it will receive a
lower impurity improvement score and be less likely to be selected.

The closest related work is the \emph{honest estimation} framework of
Athey and Imbens~\cite{athey2016}, which splits data into a ``splitting sample'' and an
``estimation sample.'' However, honest trees use the estimation sample only
\emph{after} tree construction to re-estimate leaf parameters---they do not
influence split selection. Our approach is fundamentally different: evaluation
targets directly affect which split is chosen.

% =========================================================================
\section{Method}
\label{sec:method}

\subsection{Notation}

Let $\mathcal{D} = \{(\mathbf{x}_i, y_i)\}_{i=1}^{N}$ be the training set
with feature vectors $\mathbf{x}_i \in \mathbb{R}^d$ and targets $y_i$.
A boolean mask $\mathbf{m} \in \{0,1\}^N$ partitions $\mathcal{D}$ into:
\begin{align}
    \mathcal{T} &= \{i : m_i = 1\} \quad \text{(training set)} \\
    \mathcal{E} &= \{i : m_i = 0\} \quad \text{(evaluation set)}
\end{align}

\subsection{Modified Split Selection}

In standard JPT learning, the impurity improvement for a candidate split at
position $s$ on feature $j$ is:
\begin{equation}
    \Delta I(s, j) = I(\mathcal{S}) -
    \frac{|\mathcal{S}_L|}{|\mathcal{S}|} I(\mathcal{S}_L) -
    \frac{|\mathcal{S}_R|}{|\mathcal{S}|} I(\mathcal{S}_R)
\end{equation}
where $\mathcal{S}$ is the current node's sample set and
$\mathcal{S}_L, \mathcal{S}_R$ are the left and right partitions.

With split validation, two modifications are applied:

\paragraph{Candidate restriction.}
The set of candidate split positions is restricted to
$\{s : s \in \mathcal{T}\}$. Feature values from evaluation samples are skipped
during the sorted scan but still contribute to the running statistics.

\paragraph{Target scope.}
The impurity $I(\cdot)$ is computed over a configurable subset of targets,
controlled by a \emph{validation mode}:
\begin{itemize}
    \item \textbf{Both} (default): $I$ is computed over all targets
          ($\mathcal{T} \cup \mathcal{E}$).
    \item \textbf{Training}: $I$ is computed over training targets only
          ($\mathcal{T}$).
    \item \textbf{Evaluation}: $I$ is computed over evaluation targets only
          ($\mathcal{E}$).
\end{itemize}

\subsection{Algorithm}

Algorithm~\ref{alg:split_val} shows the modified inner loop of the
core split-finding routine.

\begin{algorithm}[h]
\caption{Split validation in the impurity inner loop}
\label{alg:split_val}
\begin{algorithmic}[1]
\REQUIRE Sorted samples $\{(\mathbf{x}_{\pi(k)}, y_{\pi(k)})\}_{k=1}^{n}$,
         mask $\mathbf{m}$, mode $\in \{\text{both}, \text{train}, \text{eval}\}$
\STATE Initialize running statistics $S_L \leftarrow \emptyset$,
       $S_R \leftarrow S_{\text{total}}(\text{mode})$
\FOR{$k = 1$ \TO $n$}
    \STATE $i \leftarrow \pi(k)$ \COMMENT{sample index in sorted order}
    \IF{sample $i$ is relevant under mode}
        \STATE Update $S_L$, $S_R$ with target $y_i$
    \ENDIF
    \IF{$m_i = 1$ \AND split is valid}
        \STATE Compute $\Delta I$ from $S_L$, $S_R$
        \STATE Update best split if $\Delta I$ is improved
    \ENDIF
\ENDFOR
\end{algorithmic}
\end{algorithm}

\subsection{Implementation}

The implementation modifies the Cython-level \texttt{Impurity} class in the JPT
codebase. When the mask is \texttt{None}, the code path is identical to the
original---zero overhead. The mask is stored as a \texttt{uint8} memory view,
accessed in a \texttt{nogil} context. The \texttt{min\_samples\_leaf} constraint
uses \emph{total} sample counts (not just training), since the resulting leaf
contains all samples.

% =========================================================================
\section{Experimental Setup}
\label{sec:setup}

\subsection{Synthetic Datasets}

\paragraph{Classification.}
$N_{\text{train}} = 200$, $N_{\text{test}} = 500$, $d = 5$ features drawn from
$\mathcal{N}(0, 1)$. True boundary: $x_0 + x_1 > 0$. Label noise rate $\eta$
applied by flipping labels uniformly at random. Repeated $R = 10$ times.

\paragraph{Regression.}
$N_{\text{train}} = 200$, $N_{\text{test}} = 500$, $d = 3$ features. Target:
piecewise-linear with additive Gaussian noise $\mathcal{N}(0, \sigma)$.

\subsection{Real-World Datasets}

We use three standard UCI datasets from scikit-learn:
\begin{itemize}
    \item \textbf{Iris}: 150 samples, 4 features, 3 classes.
    \item \textbf{Wine}: 178 samples, 13 features, 3 classes.
    \item \textbf{Breast Cancer}: 569 samples, 30 features, 2 classes.
\end{itemize}
Evaluated with stratified 5-fold cross-validation.

\subsection{Evaluation Metrics}

\begin{enumerate}
    \item \textbf{Accuracy} (classification) / \textbf{MSE} (regression):
          Leaf-based prediction using the highest-prior matching leaf.
    \item \textbf{Mean log-likelihood}: $\frac{1}{N}\sum_{i=1}^{N} \log p(\mathbf{x}_i, y_i)$,
          computed via the JPT's \texttt{likelihood()} method. This metric
          evaluates the quality of the learned joint density, which is the
          native objective of JPTs.
    \item \textbf{Tree complexity}: Number of leaves.
\end{enumerate}

Statistical significance is assessed via paired $t$-tests.

% =========================================================================
\section{Results}
\label{sec:results}

\subsection{Experiment 1: Overfitting Reduction}

Table~\ref{tab:exp1} compares standard JPT learning with split validation
on noisy classification data ($\eta = 0.15$).

\begin{table}[h]
\centering
\caption{Overfitting reduction ($\eta = 0.15$, $R = 10$). Test log-likelihood
         improvement is statistically significant.}
\label{tab:exp1}
\begin{tabular}{lcccc}
\toprule
Method & Train Acc & Test Acc & Test Log-Lik & Leaves \\
\midrule
Baseline   & $.920 \pm .011$ & $.694 \pm .032$ & $-649.6 \pm 9.7$ & $59.5 \pm 1.6$ \\
Split Val. & $.915 \pm .012$ & $.706 \pm .031$ & $\mathbf{-641.7 \pm 8.0}$ & $\mathbf{55.9 \pm 1.8}$ \\
\midrule
\multicolumn{2}{l}{Paired $t$-test $p$} & $0.271$ & $\mathbf{0.023}$ & $\mathbf{<0.001}$ \\
\bottomrule
\end{tabular}
\end{table}

While the accuracy improvement is not significant at $\alpha=0.05$, the test
log-likelihood is \emph{significantly} improved ($p = 0.023$), indicating that
the learned density is a better model of the data-generating process. The
leaf count reduction is highly significant ($p < 0.001$).

\subsection{Experiment 2: Effect of Mask Ratio}

Table~\ref{tab:exp2} shows the effect of varying the training ratio.

\begin{table}[h]
\centering
\caption{Effect of mask ratio on test metrics.}
\label{tab:exp2}
\begin{tabular}{cccc}
\toprule
Ratio & Test Acc & Test Log-Lik & Leaves \\
\midrule
0.3 & $0.705 \pm 0.023$ & $\mathbf{-618.4 \pm 10.6}$ & $\mathbf{49.7 \pm 1.8}$ \\
0.5 & $0.704 \pm 0.035$ & $-636.4 \pm 9.2$ & $54.5 \pm 1.6$ \\
0.7 & $0.704 \pm 0.029$ & $-645.0 \pm 5.9$ & $56.5 \pm 1.4$ \\
0.9 & $0.700 \pm 0.033$ & $-645.9 \pm 13.3$ & $58.1 \pm 1.6$ \\
1.0  & $0.694 \pm 0.032$ & $-649.6 \pm 9.7$ & $59.5 \pm 1.6$ \\
\bottomrule
\end{tabular}
\end{table}

All three metrics show monotonic improvement as the training ratio decreases:
accuracy increases, log-likelihood improves, and leaf count decreases. The
log-likelihood improvement from ratio 1.0 to 0.3 is 31.2 nats---a substantial
improvement in density estimation quality.

\subsection{Experiment 3: Validation Mode Comparison}

\begin{table}[h]
\centering
\caption{Validation mode comparison (mask ratio $= 0.7$, $\eta = 0.15$).}
\label{tab:exp3}
\begin{tabular}{lccc}
\toprule
Mode & Test Acc & Test Log-Lik & Leaves \\
\midrule
No mask     & $0.694 \pm 0.032$ & $-649.6 \pm 9.7$ & $59.5$ \\
Both        & $0.710 \pm 0.033$ & $\mathbf{-640.3 \pm 5.5}$ & $\mathbf{56.1}$ \\
Training    & $\mathbf{0.725 \pm 0.041}$ & $-644.5 \pm 12.1$ & $57.2$ \\
Evaluation  & $0.703 \pm 0.033$ & $-645.3 \pm 10.9$ & $57.6$ \\
\bottomrule
\end{tabular}
\end{table}

The \emph{training} mode achieves the highest accuracy (0.725), while the
\emph{both} mode achieves the best log-likelihood ($-640.3$). This suggests
that different modes optimize different objectives: the \emph{both} mode
produces a better density model, while the \emph{training} mode produces a
better classifier.

\subsection{Experiment 4: Noise Level Sensitivity}

\begin{table}[h]
\centering
\caption{Test metrics across noise levels (mask ratio $= 0.7$, mode $=$ both).}
\label{tab:exp4}
\begin{tabular}{ccccc}
\toprule
$\eta$ & Base Acc & SV Acc & Base Log-Lik & SV Log-Lik \\
\midrule
0.00 & $0.926$ & $0.920$ & $-658.5$ & $\mathbf{-642.0}$ \\
0.05 & $0.852$ & $0.852$ & $-653.3$ & $\mathbf{-644.3}$ \\
0.10 & $0.748$ & $\mathbf{0.769}$ & $-649.3$ & $\mathbf{-642.7}$ \\
0.15 & $0.694$ & $0.697$ & $-649.6$ & $\mathbf{-642.4}$ \\
0.20 & $0.660$ & $0.654$ & $-645.4$ & $\mathbf{-640.6}$ \\
0.25 & $0.612$ & $0.616$ & $-651.1$ & $\mathbf{-642.7}$ \\
0.30 & $0.576$ & $0.574$ & $-654.1$ & $\mathbf{-645.6}$ \\
\bottomrule
\end{tabular}
\end{table}

The log-likelihood improvement from split validation is \textbf{consistent
across all noise levels}, averaging 8.4 nats. This robustness suggests
that the density improvement is a structural property of the regularization,
not an artifact of a specific noise regime. Accuracy improvements are
concentrated at moderate noise ($\eta = 0.10$: $+2.1$ pp).

\subsection{Experiment 5: Regression}

\begin{table}[h]
\centering
\caption{Regression results (mask ratio $= 0.7$, mode $=$ both).}
\label{tab:exp5}
\begin{tabular}{ccccc}
\toprule
$\sigma$ & Base MSE & SV MSE & Base Log-Lik & SV Log-Lik \\
\midrule
0.5 & $0.799$ & $0.775$ & $-508.5$ & $\mathbf{-490.5}$ \\
1.0 & $1.744$ & $1.738$ & $-503.3$ & $\mathbf{-487.4}$ \\
2.0 & $6.005$ & $5.827$ & $-507.7$ & $\mathbf{-496.5}$ \\
\bottomrule
\end{tabular}
\end{table}

Log-likelihood improvements are substantial (12--18 nats), with the largest
improvement at low noise where the density model is most informative.
MSE improvements are modest but consistent.

\subsection{Experiment 6: Real-World Datasets}

\begin{table}[h]
\centering
\caption{Real-world datasets (stratified 5-fold CV, mask ratio $= 0.7$).}
\label{tab:exp6}
\begin{tabular}{llcccc}
\toprule
Dataset & Method & Acc & Log-Lik & Leaves \\
\midrule
\multirow{2}{*}{Iris} & Baseline & $.973 \pm .025$ & $-449.9$ & $33.4$ \\
 & Split Val. & $.947 \pm .040$ & $-445.2$ & $33.0$ \\
\midrule
\multirow{2}{*}{Wine} & Baseline & $.877 \pm .038$ & $-690.8$ & $43.6$ \\
 & Split Val. & $.860 \pm .030$ & $-686.9$ & $41.0$ \\
\midrule
\multirow{2}{*}{Breast Cancer} & Baseline & $.903 \pm .026$ & $-635.1$ & $46.6$ \\
 & Split Val. & $.903 \pm .026$ & $-628.8$ & $44.6$ \\
\bottomrule
\end{tabular}
\end{table}

On these datasets, split validation consistently improves log-likelihood and
reduces tree complexity, though the accuracy effect is mixed. The Breast Cancer
dataset shows the clearest benefit: identical accuracy with improved density and
fewer leaves.

\subsection{Experiment 7: Mode Comparison on Breast Cancer}

\begin{table}[h]
\centering
\caption{Mode comparison on Breast Cancer (5-fold CV, mask ratio $= 0.7$).}
\label{tab:exp7}
\begin{tabular}{lccc}
\toprule
Mode & Accuracy & Log-Lik & Leaves \\
\midrule
No mask     & $0.903 \pm 0.026$ & $-635.1$ & $46.6$ \\
Both        & $0.907 \pm 0.033$ & $\mathbf{-627.8}$ & $\mathbf{44.8}$ \\
Training    & $\mathbf{0.916 \pm 0.035}$ & $-622.9$ & $45.8$ \\
Evaluation  & $\mathbf{0.916 \pm 0.028}$ & $-631.3$ & $45.6$ \\
\bottomrule
\end{tabular}
\end{table}

On real data, both the \emph{training} and \emph{evaluation} modes achieve 1.3
percentage points higher accuracy than baseline, with the \emph{training} mode
also delivering the best log-likelihood ($-622.9$ vs.\ $-635.1$).

\subsection{Experiment 8: Mask Ratio on Wine}

\begin{table}[h]
\centering
\caption{Mask ratio on Wine (5-fold CV, mode $=$ both). The strongest
         regularization ratio (0.3) achieves 6.2 pp higher accuracy than
         baseline.}
\label{tab:exp8}
\begin{tabular}{ccc}
\toprule
Ratio & Accuracy & Leaves \\
\midrule
0.3 & $\mathbf{0.939 \pm 0.032}$ & $\mathbf{36.8}$ \\
0.5 & $0.916 \pm 0.050$ & $39.8$ \\
0.7 & $0.888 \pm 0.047$ & $40.6$ \\
0.9 & $0.877 \pm 0.038$ & $42.8$ \\
1.0 (no mask) & $0.877 \pm 0.038$ & $43.6$ \\
\bottomrule
\end{tabular}
\end{table}

This is the most striking result: on the Wine dataset, a training ratio of 0.3
improves accuracy from 0.877 to \textbf{0.939}---a 6.2 percentage point gain,
while reducing leaves by 15.6\%. The pattern of lower ratio $\to$ better
accuracy is monotonic, suggesting that the Wine dataset's 13 features contain
substantial noise that benefits from aggressive regularization.

% =========================================================================
\section{Discussion}
\label{sec:discussion}

\subsection{Key Findings}

\begin{enumerate}
    \item \textbf{Log-likelihood is the most reliable improvement metric.}
          Across synthetic and real data, split validation consistently
          improves test log-likelihood, often significantly ($p = 0.023$ in
          Exp.~1). This confirms that the method produces a better density
          model, which is the native objective of JPT learning.

    \item \textbf{Structural regularization is highly significant.}
          The leaf count reduction ($p < 0.001$) is the strongest statistical
          signal. Simpler trees are more interpretable and less prone to
          overfitting.

    \item \textbf{Real-world accuracy gains can be substantial.}
          The Wine dataset shows a 6.2~pp accuracy improvement at ratio 0.3,
          demonstrating that the method can make a practical difference on
          real data, particularly on datasets with many noisy features relative
          to samples.

    \item \textbf{The mask ratio is a smooth regularization knob.}
          Lower ratios consistently produce simpler, often more accurate models.
          This continuous control is preferable to discrete hyperparameters.

    \item \textbf{Different modes suit different objectives.}
          The \emph{training} mode maximizes classification accuracy, while
          the \emph{both} mode maximizes log-likelihood. The \emph{evaluation}
          mode is competitive on real data but less consistent on synthetic data.
\end{enumerate}

\subsection{Why Log-Likelihood Improves More Than Accuracy}

Accuracy measures only the most probable class, discarding calibration
information. Log-likelihood captures the full predictive distribution: a tree
with well-calibrated class probabilities (e.g., 0.8/0.2 for a true 80/20 split)
scores better than one with overconfident predictions (1.0/0.0). Split
validation's regularization effect produces better-calibrated leaf distributions,
which improves log-likelihood even when the argmax prediction is unchanged.

\subsection{Comparison to Related Work}

\paragraph{Honest trees.}
Athey and Imbens~\cite{athey2016} split data for post-hoc estimation.
Split validation influences tree \emph{structure} during construction---a
fundamentally different and complementary mechanism.

\paragraph{Pre-pruning.}
Standard pre-pruning applies uniform constraints. Split validation provides
data-adaptive regularization: noisy regions produce fewer candidate splits,
while clean regions retain full expressiveness.

\paragraph{Random forests.}
Random forests achieve regularization via ensemble averaging. Split validation
is complementary---it could be applied within individual trees of an ensemble.

\subsection{Limitations}

\begin{itemize}
    \item Real-world evaluation is limited to three UCI datasets. Larger and
          more complex datasets would strengthen the conclusions.
    \item The optimal mask ratio depends on the dataset's noise-to-signal
          ratio. An adaptive selection strategy remains an open question.
    \item With 5-fold CV, some real-world comparisons lack statistical power.
    \item The current implementation uses a fixed mask. Depth-dependent or
          bootstrap-based strategies remain unexplored.
\end{itemize}

% =========================================================================
\section{Conclusion}
\label{sec:conclusion}

Split validation is a lightweight, principled regularization technique for
tree-based density estimation. By restricting candidate split points to a
training subset while evaluating impurity over all targets, it creates an
inductive bias toward generalizable splits. Our experiments demonstrate that
the method:
\begin{itemize}
    \item Significantly improves test log-likelihood ($p = 0.023$), confirming
          better density estimation.
    \item Reliably reduces tree complexity ($p < 0.001$).
    \item Achieves substantial accuracy gains on real data (up to $+6.2$~pp on
          Wine).
    \item Provides a smooth regularization knob via the mask ratio.
\end{itemize}

The method is easy to implement, introduces zero overhead when disabled, and
is complementary to existing regularization techniques. Future work should
investigate adaptive masking strategies, interaction with ensemble methods, and
evaluation on large-scale real-world datasets.

% =========================================================================
\bibliographystyle{plain}
\begin{thebibliography}{9}

\bibitem{athey2016}
S.~Athey and G.~Imbens.
\newblock Recursive partitioning for heterogeneous causal effects.
\newblock {\em Proceedings of the National Academy of Sciences},
  113(27):7353--7360, 2016.

\bibitem{breiman1984}
L.~Breiman, J.~Friedman, C.~J.~Stone, and R.~A.~Olshen.
\newblock {\em Classification and Regression Trees}.
\newblock CRC Press, 1984.

\bibitem{jpt2022}
T.~Nickl, M.~Kl\"as, and A.~Pretschner.
\newblock Joint probability trees.
\newblock {\em Machine Learning}, 2022.

\end{thebibliography}

\end{document}
